\documentclass[tikz,border=4pt]{standalone}
\usepackage{ctex}
\usepackage{amsmath}
\usetikzlibrary{shapes.geometric, arrows.meta, positioning}

\tikzset{
  box/.style={rectangle, draw=black, rounded corners=3pt,
              minimum width=3.8cm, minimum height=0.65cm,
              align=center, font=\small},
  arr/.style={->, >=Stealth, thick},
}

\begin{document}
% 代入消元法 vs 加减消元法 对比流程
\begin{tikzpicture}[node distance=0.9cm]

  \node[box, fill=blue!10, minimum width=7cm] (start)
    {二元一次方程组 $\begin{cases}a_1 x+b_1 y=c_1\\a_2 x+b_2 y=c_2\end{cases}$};

  \node[box, fill=yellow!15, below=of start, minimum width=7cm] (obs)
    {观察：某方程中某未知数系数是 $\pm 1$？};

  % 两个分支
  \node[box, fill=green!15, below left=1.0cm and 1.5cm of obs] (sub)
    {\textbf{代入消元法}\\从系数 $\pm1$ 方程\\解出 $y=f(x)$，代入另一式};

  \node[box, fill=red!10, below right=1.0cm and 1.5cm of obs] (elim)
    {\textbf{加减消元法}\\乘适当倍数使\\某未知数系数相同/相反\\再相加减消元};

  \node[box, fill=blue!10, minimum width=7cm,
        below=3.0cm of obs] (end)
    {得到一元方程，解之；\\代回求另一未知数；验算};

  \draw[arr] (start) -- (obs);
  \draw[arr] (obs) -- node[above left, font=\small]{是} (sub);
  \draw[arr] (obs) -- node[above right, font=\small]{否} (elim);
  \draw[arr] (sub) -- (end);
  \draw[arr] (elim) -- (end);

\end{tikzpicture}
\end{document}
